% q3_efg.tex - Question 3 Studies E-G (OWNER: C) - revised from docs/C_Q3_EFG.md
\section{Question 3, Studies E--G}

All tests use $\alpha = 0.05$. Each study uses the complete set of records available in the named dataset for its analysis variables.

\FloatBarrier
\subsection*{Study E --- Representation of public and private colleges}

\texttt{Private} is a nominal variable with two categories. Each college is treated as an independent observation. Using all 777 colleges in \texttt{ISLR2::College}, a chi-square goodness-of-fit test compares the observed counts with a specified 50:50 split. The null hypothesis ($\Hnull$) assigns a probability of 0.5 to each category. The alternative ($\Halt$) allows the probabilities to differ from 0.5.

\begin{table}[htbp]
  \centering
  \caption*{Table Q3-E. Observed and expected college counts}
  \begin{tabular}{lrr}
    \toprule
    College type & Observed & Expected under $\Hnull$ \\
    \midrule
    Public & 212 & 388.5 \\
    Private & 565 & 388.5 \\
    \bottomrule
  \end{tabular}
\end{table}

Both expected counts exceed 5. The test gave $\chisq(1) = 160.372$, $\pval < 0.001$. We reject $\Hnull$. The data provide strong evidence against a 50:50 split. Private colleges account for 72.72\% of the records.

\FloatBarrier
\subsection*{Study F --- College status and incoming student profile}

\texttt{Private} and the High/Lower student profile are categorical variables. A college is classified as High if at least 50\% of its incoming students were in the top 10\% of their high school class. Using the same 777 colleges, we apply Pearson's chi-square test of independence without Yates' correction, treating colleges as independent observations. $\Hnull$ states that college status and student profile are independent. $\Halt$ states that they are associated.

\begin{table}[htbp]
  \centering
  \caption*{Table Q3-F. Student profile by college type}
  \begin{tabular}{lrrr}
    \toprule
    College type & Lower & High & Percent High \\
    \midrule
    Public & 198 & 14 & 6.60\% \\
    Private & 496 & 69 & 12.21\% \\
    \bottomrule
  \end{tabular}
\end{table}

The minimum expected count was 22.65. The test gave $\chisq(1) = 5.083$, $\pval = 0.024$. We reject $\Hnull$. The proportion of colleges classified as High was 12.21\% among private colleges and 6.60\% among public colleges. Study E tests whether a single variable follows a specified distribution, while Study F tests for an association between two variables. Because this is an observational comparison, the difference cannot be attributed to private status alone. Admissions selectivity, for example, may also contribute.

\FloatBarrier
\subsection*{Study G --- Variability in fuel economy by vehicle origin}

Fuel economy (\texttt{mpg}) is a quantitative ratio variable. Vehicle origin is nominal, with three groups treated as independent. Using all 392 vehicles in \texttt{ISLR2::Auto}, we apply a median-centred Levene test, which compares variability across groups without requiring normally distributed mpg. $\Hnull$ states that the three population variances are equal. $\Halt$ states that at least one differs.

\begin{table}[htbp]
  \centering
  \caption*{Table Q3-G. Fuel economy variance by origin}
  \begin{tabular}{lrr}
    \toprule
    Origin & Sample size & Sample variance (mpg$^2$) \\
    \midrule
    American & 245 & 41.479 \\
    European & 68 & 43.299 \\
    Japanese & 79 & 37.089 \\
    \bottomrule
  \end{tabular}
\end{table}

The test gave $F(2, 389) = 0.086$, $\pval = 0.918$. We fail to reject $\Hnull$. Although the European group had the largest sample variance, the test provides insufficient evidence of a difference in population variances across the three origins.
